% SPDX-FileCopyrightText: 2026 Will Cook
% SPDX-License-Identifier: CC-BY-4.0
%
% Working record behind the synthesis note.  This is a reasoning surface: a
% record of attempted routes, the routes that closed, and what survives.
\documentclass[11pt]{article}
% SPDX-FileCopyrightText: 2026 Will Cook
% SPDX-License-Identifier: CC-BY-4.0
%
% Shared preamble for the Erdős Problem Notes: the short, standalone,
% problem-owned manuscripts covering the expansion library `ErdosProblems`.
%
% One note owns one Erdős problem.  Each note repeats whatever context its own
% reader needs, so that a reader who arrives at a single problem never has to
% assemble it from the gateway paper.  Deliberate repetition across notes is the
% design, not an oversight.
%
% Source links resolve against one pinned formal-source commit, declared once
% below.  `scripts/check_problem_note_sources.py` verifies that every authored
% (file, line, declaration) triple names that exact declaration in the pinned
% snapshot, so a link cannot silently rot as later work moves lines.

\usepackage{paper-house-style}
\usepackage{array}
\usepackage{booktabs}
\usepackage{enumitem}
\setlist{topsep=0.35em,itemsep=0.2em}
\usepackage[colorlinks=true,linkcolor=linkinternal,urlcolor=linkexternal,
  citecolor=linkinternal]{hyperref}
\hypersetup{bookmarksnumbered=true,bookmarksopen=true,bookmarksopenlevel=1,
  pdfauthor={Will Cook},pdflang={en-GB}}

% ---------------------------------------------------------------------------
% Pinned formal source.  \commit names the checkpoint every link resolves
% against; the checker reads that snapshot rather than the working tree, so the
% printed coordinates stay correct however later waves move declarations.
% ---------------------------------------------------------------------------
\newcommand{\commit}{e5fd26a6d1b1a346430205eab5385b4c9565d004}
\newcommand{\commitshort}{e5fd26a6d1b1}
\newcommand{\repobase}{https://github.com/wcook04/plectis-lean-erdos249-257/blob/\commit}
\newcommand{\sourceurl}{https://github.com/wcook04/plectis-lean-erdos249-257/tree/\commit}
\newcommand{\PX}{ErdosProblems}

% Declaration names stay inside link targets.  The printed label is either a
% spaced, lower-case reading of the declaration name (\lref) or a short authored
% phrase (\lword); neither prints a module path or a raw Lean identifier.
\ExplSyntaxOn
\cs_new_protected:Npn \noteleanlabel #1
  {
    \tl_set:Nx \l_tmpa_tl { \tl_to_str:n {#1} }
    \regex_replace_all:nnN { _+ } { \c{space} } \l_tmpa_tl
    \regex_replace_all:nnN
      { ([a-z0-9])([A-Z]) }
      { \1\c{space}\2 }
      \l_tmpa_tl
    \tl_set:Nx \l_tmpa_tl { \tl_lower_case:n { \tl_use:N \l_tmpa_tl } }
    \tl_use:N \l_tmpa_tl
  }
\ExplSyntaxOff
\newcommand{\leanlabel}[1]{\noteleanlabel{#1}}
\newcommand{\lref}[3]{\href{\repobase/\PX/#1\#L#2}{\leanlabel{#3}}}
\newcommand{\lword}[4]{\href{\repobase/\PX/#1\#L#2}{#4}}
\newcommand{\lloc}[2]{\href{\repobase/\PX/#1\#L#2}{Lean source}}

% Links into a sibling library, named repository-relative.  The expansion
% library is implicit in \lref/\lword, because that is where problem-owned
% work lands.  Several problems have their headline theorem in the older
% reviewed #249/#257 corpus, and a note that could not name that library
% would have to paraphrase a checked statement instead of linking it.
\newcommand{\mref}[3]{\href{\repobase/#1\#L#2}{\leanlabel{#3}}}
\newcommand{\mword}[4]{\href{\repobase/#1\#L#2}{#4}}
\newcommand{\mloc}[2]{\href{\repobase/#1\#L#2}{Lean source}}
\emergencystretch=2em

\newtheorem{theorem}{Theorem}[section]
\newtheorem{proposition}[theorem]{Proposition}
\newtheorem{lemma}[theorem]{Lemma}
\newtheorem{corollary}[theorem]{Corollary}
\theoremstyle{definition}
\newtheorem{definition}[theorem]{Definition}
\newtheorem{example}[theorem]{Example}
\newtheorem{problem}[theorem]{Problem}
\theoremstyle{remark}
\newtheorem*{remark}{Remark}

\newcommand{\N}{\mathbb{N}}
\newcommand{\Npos}{\mathbb{N}_{>0}}
\newcommand{\Z}{\mathbb{Z}}
\newcommand{\Q}{\mathbb{Q}}
\newcommand{\R}{\mathbb{R}}
\newcommand{\lcm}{\operatorname{lcm}}
\newcommand{\ph}{\varphi}

% The series line printed under every note's title block.
\newcommand{\seriesline}{Erd\H{o}s Problem Notes}

% Production provenance.  The wording is fixed by the corpus provenance
% contract and reproduced verbatim in every manuscript in this repository, so
% the papers cannot drift into different accounts of how they were made.
\newcommand{\noteprovenance}{\provenancenote{\emph{Authorship and AI use.} The
prose, formal proofs, and software in this version were drafted and revised by
large-language-model agents under Will Cook's direction.  Cook set the
objectives and acceptance criteria, reviewed and authorised the manuscript, and
is responsible for its contents.  The agents are production tools, not authors.
See \emph{Statements and declarations}.}}

% The status paragraph every note carries, in its own words, before any result.
% Kernel checking establishes the exact Lean proposition; it does not establish
% that the proposition is the interesting one, that it is new, or that the
% Erdős problem is closed.
\newcommand{\statusboundary}{%
  \emph{Status.}  The problem treated here is open, and this note does not
  close it.  Every statement below marked as checked is a proposition that the
  pinned Lean kernel accepts from the sources this note links to, with no
  \texttt{sorry}, no added axiom, and no unchecked evaluation.  That is a claim
  about the formal statement, not about its mathematical interest, its
  novelty, or the original problem.  The unresolved obligations are named
  exactly, in their own section, and none of the finite computations,
  reductions, or no-go results here removes one of them.}

\usepackage{xurl}
\usepackage{longtable}
\hypersetup{
  pdftitle={Reading eight Erdős problems together: a working record},
  pdfsubject={Erdős Problem Notes: record behind the synthesis note},
  pdfkeywords={irrationality, subsum sets, negative results, prior art, Farey fractions},
}

% Immutable formal-source revision used by the corpus; this note links no Lean
% source itself and cites sibling papers by file name.
\renewcommand{\commit}{99f4bf47422abbd8757cbb22b50ba079d764d3a7}
\renewcommand{\commitshort}{99f4bf47422a}

\runninghead{Reading eight Erd\H{o}s problems together: record}
\title{Reading Eight Erd\H{o}s Problems Together}
\subtitle{A working record: proofs, eliminations, rediscoveries and open questions}
\author{Will Cook}
\date{20 September 2026}

\begin{document}
\maketitle
\noteprovenance

\begin{abstract}
This record supports the note \cite{plectisshort}.  It proves the two lemmas
about base $2$ that the note uses without proof, lists the negative results of
the seven companion notes with their locations and status, gives in full the
two wrong readings of one computation and the theorem that replaced them,
records the ideas that were eliminated together with the statement that
eliminated each, records what turned out to be known already, and ends with the
questions we think are worth someone's time.  It is a lookup document.  Each
fact has one home here.
\end{abstract}

\section{What the note relies on}\label{sec:map}

\begin{center}
\small
\begin{tabular}{p{0.34\linewidth}p{0.58\linewidth}}
\hline
In the note & Here \\
\hline
Theorem 1.1 (choices against contraction) & Proved in the note.  Attribution in Section~\ref{sec:known}. \\
Theorem 1.2(b), base $2$ & Quoted from \cite[Section~7]{plectis257}.  Removed intervals: Lemma~\ref{lem:gaps}. \\
Theorem 1.3 (chains) & Proved in the note.  A wrong citation caught on the way: Section~\ref{sec:eliminated}, item 4. \\
Theorem 1.4 (forced statistic) & Proved in the note.  The readings it corrects: Section~\ref{sec:misread}. \\
Limits on methods & Section~\ref{sec:inventory}, with status and sources. \\
\hline
\end{tabular}
\end{center}

Every proof here is an ordinary proof.  Nothing is checked in Lean.  The
computations are exact rational arithmetic with saved outputs
\cite{plectisinvestigation}.

\section{Two lemmas at base two}\label{sec:lemmas}

Notation is that of the note at $t=2$: $w_n=(2^n-1)^{-1}$,
$R_N=\sum_{n>N}w_n$, $g_n=w_n-R_n>0$, $X_F=\sum_{n\in F}w_n$ for finite $F$,
and $\mathcal A$ is the set of all subsums.

\begin{lemma}[the removed intervals]\label{lem:gaps}
$[0,E]\setminus\mathcal A$ is the disjoint union, over finite nonempty $F$, of
the open intervals $(X_F-g_{\max F},\,X_F)$.  Their total length is
$\sum_n2^{n-1}g_n=E-1$, and
$g_n=\sum_{j\ge2}\frac{2^j-2}{2^j-1}2^{-jn}=\tfrac23 4^{-n}+\tfrac67 8^{-n}+\cdots$.
\end{lemma}

\begin{proof}
Fix the digits $\varepsilon_1,\dots,\varepsilon_{n-1}$ and let
$s=\sum_{i<n}\varepsilon_iw_i$.
The values with these digits lie in $[s,s+R_{n-1}]$, and they split into
$[s,s+R_n]$ and $[s+w_n,s+w_n+R_n]$.  The interval between them is
$(s+R_n,s+w_n)$.  Its right end is $X_F$ with
$F=\{i<n:\varepsilon_i=1\}\cup\{n\}$ and its length is $g_n$.  Every finite
nonempty $F$ arises once.  The total length is $E$ minus the measure of
$\mathcal A$, which is $1$ \cite[Section~7]{plectis257}.  The series for $g_n$
follows from $w_n=\sum_{j\ge1}2^{-jn}$ and
$R_n=\sum_{j\ge1}2^{-jn}/(2^j-1)$.
\end{proof}

\begin{lemma}[translation by a finite subsum]\label{lem:translate}
Let $F$ be finite with largest element $n$ and let $0\le x\le R_n$.  The greedy
rule applied to $X_F+x$ selects exactly $F$ among the indices up to $n$ and then
agrees with the greedy rule applied to $x$.  In particular $X_F+x\in\mathcal A$
if and only if $x\in\mathcal A$, and $X_F+x$ is rejected at a step $m>n$ exactly
when $x$ is.
\end{lemma}

\begin{proof}
At an index $i\le n$ the remainder is $X_{F\cap[i,n]}+x$.  If $i\in F$ it is at
least $w_i$ and the index is taken.  If $i\notin F$ it is at most
$R_i-R_n+x\le R_i<w_i$, so the index is skipped and nothing is rejected.  After
index $n$ the remainder is $x$.
\end{proof}

Two consequences.  The map $x\mapsto x+1$ preserves reduced denominators, so at
every step $n\ge2$ the number of fractions of height at most $Q$ rejected at
that step is even; every saved table has this property.  The condition is
$x\le R_{\max F}$.  It is not enough that $X_F+x\le E$: $x=1/2$ is not rejected
through step $160$, while $1/2+1/3=5/6$ lies in $(R_1,1)$ and is rejected at
step $1$.  Counts of rejections should therefore be taken modulo these
translations.  At height $200$ the four fractions rejected at step $12$ are
$46/183$ and its translates by $1/3$, $1$ and $4/3$, one event.

A finite subsum has odd denominator, since $\prod_{n\in F}(2^n-1)$ is odd.  So a
rational with even denominator is never a finite subsum, and if it is a subsum
at all its set $S$ is infinite.

\section{Negative results in the companion notes}\label{sec:inventory}

Status is as stated in each note: L for checked in Lean there, O for an
ordinary proof there, C for cited there.

\begin{longtable}{p{0.10\linewidth}p{0.20\linewidth}p{0.52\linewidth}p{0.06\linewidth}}
\hline
Note & Location & Statement & \\
\hline
\#68 \cite{plectis68} & Section 6 & The integer-gap comparison cannot hold at cutoffs within a bounded distance of the degree.  Small tails alone do not give the strict comparison. & O \\
\#243 \cite{plectis243} & Proposition 17 & A counterexample has errors that are eventually nonzero, relatively small, with unbounded negative parts.  Size of the error does not decide; the denominators must be used. & L \\
\#249 \cite{plectis249} & Section 5 & A rational series with the totient's values at odd indices, within $2$ at even indices, and sum $5/4$.  Positive tail differences need not be nonintegral. & L \\
\#249 \cite{plectis249} & Theorem 9 & Every admissible rank-one quotient stays more than $21/320$ from its target. & L \\
\#251 \cite{plectis251} & Proposition 1, Corollary 2 & Sparse nonnegative corrections below any unbounded allowance reach every real in an interval and keep all congruences.  A bounded allowance is impossible.  Sources: \cite{fridy1966}, \cite[Lemma~4]{crmarickovac2025}, \cite[Lemma~5.1]{kovactao2024}, \cite{vandoornkovac2025}. & O \\
\#257 \cite{plectis257} & Section 6 & The small-displacement quantity stays above $1/2$ at full support, where the value is irrational \cite{erdos1948}.  No proof covering full support can rest on it. & L, O \\
\#257 \cite{plectis257} & Section 3 & Every positive divisor cover costs at least $e$ times the mean of $\log^+$ of its multiplicity.  The averaging method cannot reach the prime support, where irrationality is known at base $2$ \cite{taoteravainen2025}. & O \\
\#269 \cite{plectis269} & Theorem 1 & Nonsingular minors of every order: no finite sum of products separates one exponent from the other two.  Fan posted the two-prime separation \cite{fan2026comment}; the three-prime statement is the note's. & L \\
\#1049 \cite{plectis1049} & Theorem 5 & A family of forms with degrees at most $\delta n^2$ and decay $-\sigma n^2\log x$ at every $x>1$ has $\sigma\le\delta$; its region never reaches $\log b/\log a=1/2$.  Ordinary proof, with the contradiction step and the comparison with $1/2$ checked in Lean. & O, L \\
\#1049 \cite{plectis1049} & Theorem 7 & The stated clearing conditions cannot be met at base $3/2$. & L \\
\hline
\end{longtable}

Rows \#249 Section 5 and \#251 are instances of Theorem 1.1(ii) of the note.
The others bound a method.  We tried to state one inequality that covers
\#1049 Theorem 5 and the cover cost of \#257, a cost of clearing denominators
against the decay gained, and did not find a formulation that survives both
sets of hypotheses.  We do not claim the rows share a cause.

\section{One computation read wrongly twice}\label{sec:misread}

The computation \cite{plectisinvestigation} runs the greedy rule in exact
arithmetic on every reduced fraction in $(0,E]$ with $2\le q\le Q$.

\emph{First reading.}  At $Q=36$, $382$ of $633$ fractions are not rejected
through step $160$ and $14$ are finite subsums, a share $(382+14)/633=0.6256$
close to $1/E=0.6224$.  This was read as evidence that about $62\%$ of rationals are
subsums, hence that \#257 is false.  The reading is wrong.  By Theorem 1.4 of
the note the share at any fixed depth tends to $2^NR_N/E$ whatever the truth of
\#257, because fractions equidistribute over the $2^N$ intervals that survive
$N$ steps.

\emph{Second reading.}  After that correction, the rejections found at steps
$8$, $9$ and $12$ once $Q$ reached $200$ were described as the informative
data, and the $12{,}218$ fractions not rejected through step $60$ as survivors
of interest.  This is also wrong.  The expected number rejected at step $n$ is
$2^{n-1}g_n\sum_{2\le q\le Q}\varphi(q)$, asymptotically
$(3Q^2/\pi^2)2^{n-1}g_n$, which is below $1$ for $n>2\log_2Q-2\log_2\pi$, about
step $12$ at $Q=200$.  Survival past that step is forced by counting, and the
rejections before it confirm the measures of Lemma~\ref{lem:gaps} and nothing
else.  The error term $O(2^N\log Q/Q)$ of Theorem 1.4 is useful at $Q=200$ only
for $N$ up to about $5$.  At $N=12$ the observed share is $0.62372$ against
$2^{12}R_{12}/E=0.62245$, an agreement better than the theorem guarantees.

What survives is the agreement itself.  At $Q=200$ the counts at steps $1$ to
$9$ are $4809$, $1470$, $600$, $268$, $132$, $66$, $32$, $8$, $6$, against
$4811$, $1467$, $604$, $277$, $133$, $65$, $32$, $16$, $8$ from the measures of
Lemma~\ref{lem:gaps}; steps $10$ and $11$ have none against $4$ and $2$, and
step $12$ has $4$ against $1$.  The late counts fluctuate more than independent
events would, because rejections arrive in the families of
Lemma~\ref{lem:translate}.

\section{Eliminated ideas}\label{sec:eliminated}

\begin{enumerate}
\item \emph{If the reachable values form a null set, no rational value is
  reachable.}  False.  Take the binary series with digit $1$ everywhere except
  digit $0$ at positions $n_1<n_2<\cdots$, and allow each of those digits to be
  changed to $1$.  The reachable values form a null set when the positions are
  sparse, by Theorem 1.1(i), yet changing all of them gives $\sum2^{-n}=1$.
  So no count of choices against contraction, and no depth depending only on
  sparsity, excludes a particular rational.
\item \emph{Every infinite subset of a host with null subsum set has an
  irrational sum.}  A host is a set $B$ of allowed indices, and its subsum set
  is $\{X_S(2):S\subseteq B\}$.  For hosts chosen without reference to the target this is
  open and is a form of \#257 itself.  As a universal statement it cannot be a
  route: the support of any rational subsum with infinite $S$ would be such a
  host.  The subsum set of a host $B$ has positive measure exactly when the
  complement of $B$ is finite, by Theorem 1.1(i) and the measure at full
  support.
\item \emph{The share of surviving fractions as evidence.}  Section~\ref{sec:misread}.
\item \emph{A wrong locator.}  A draft of Theorem 1.3(b) cited Theorem 1.1 of
  \cite{adamczewskifaverjon2026} as Nishioka's theorem.  That theorem says a
  value of a Mahler function at an algebraic point is rational or
  transcendental, which cannot prove irrationality.  The proof in the note uses
  Mahler's theorem \cite{mahler1929}, whose hypotheses are checked there.
\item \emph{Algebraic independence for \#1049.}  With $g$ as in the note,
  $\sum_{n\ge1}(t^n-1)^{-1}=\sum_{m\ \mathrm{odd}}g(t^{-m})$, and each $g(t^{-m})$
  is transcendental for rational $t>1$.  This gives nothing for the infinite
  sum: limits of transcendental numbers take every value, and no form of
  Mahler's method controls an infinite sum of Mahler values.
\end{enumerate}

\section{What was already known}\label{sec:known}

\begin{itemize}
\item The covering argument of Theorem 1.1(ii) goes back to Kakeya; see
  \cite{hornich1941,nitecki2013,fridy1966}.  Its use to build rational series
  inside a class defined by soft data is the method of Kova\v{c} and Tao
  \cite{kovactao2024}, of Crmari\'{c} and Kova\v{c} \cite{crmarickovac2025} and
  of van Doorn and Kova\v{c} \cite{vandoornkovac2025}.
\item That the subsums of $\sum(t^n-1)^{-1}$ form a Cantor set at fixed
  $t\ge2$ is \cite[Remark~4.1]{kovactao2024}.
\item A closed set of positive measure can contain essentially no rationals
  \cite{boesdarsterdos1981}, so the heuristic of the note's Section 4 cannot be
  a consequence of measure.
\item Counting rationals of bounded height is studied for null Cantor sets
  \cite{rstw2019,chowvarjuyu2024}.  We found no work on rational points of a
  subsum set of positive measure, and none on rational points in the subsum-set
  literature.  The searches were made on 20 September 2026 and are listed in
  the repository record.
\item One identity with a consequence.  Since
  $\sum_{n\ge1}\mu(n)/(b^n-1)=1/b$, the sums of $(b^n-1)^{-1}$ over squarefree
  $n$ with an even, respectively odd, number of prime factors are
  $\tfrac12(X_{\mathrm{sf}}(b)\pm1/b)$.  Duverney and Tachiya prove that
  $X_{\mathrm{sf}}(2^j)$ is irrational \cite[Corollary~1.2 and
  Example~1.1]{duverneytachiya}, as quoted in \cite[Section~6]{plectis257}, so
  both sums are irrational at every base $2^j$.  Both supports have divergent
  reciprocal sums.  The identity at base $2$ is derived in
  \cite[Section~5]{plectis249}; the corollary may be known.
\end{itemize}

\section{Questions}\label{sec:questions}

\begin{enumerate}
\item Is $1/2$ a subsum of $\sum(2^n-1)^{-1}$?  The exact obligation is in
  \cite[Theorem~7]{plectis257}.  By Lemma~\ref{lem:gaps} the general question is
  one-sided approximation of a rational by finite subsums $X_F$ to within
  $g_{\max F}$.
\item Does the count of fractions of height at most $Q$ rejected at step $n$
  stay close to $(3Q^2/\pi^2)2^{n-1}g_n$ in the joint range
  $n\le(2-\varepsilon)\log_2Q$, counted modulo the translations of
  Lemma~\ref{lem:translate}?  A persistent excess would be the first sign of an
  arithmetic mechanism for \#257.
\item Is $a^2>b^3$ necessary in Theorem 1.3(a)?  First case: base $4/3$, a chain
  with ratios $2$ and infinitely many ratios $3$.
\item Is there one inequality behind \#1049 Theorem 5 and the cover cost of
  \#257?
\item For factorial radices, is an allowance of order $n$ the exact threshold
  for a rational series within that allowance of a given one?  Theorem 1.1
  gives an interval when the allowance at level $n$ is at least $n$ and a null
  set when it is $o(n)$ along a sequence.  This would explain why \#68 and
  \#251 resist opposite methods.
\end{enumerate}

\begin{thebibliography}{99}
\bibitem{plectisshort} W.~Cook,
  \href{erdos-synthesis-subsums-across-bases.pdf}{\emph{Rational
  and irrational subsums of a Lambert series across bases}}, Erd\H{o}s Problem
  Notes, 2026.

\bibitem{plectisinvestigation} W.~Cook,
  \emph{Choices, contraction and rational membership}, exact computation with saved outputs in \texttt{research/experiments/choices\_contraction/} of the repository,
  2026.

\bibitem{erdos1948} P.~Erd\H{o}s, \emph{On arithmetical properties of Lambert
  series}, J.~Indian Math.\ Soc.\ 12 (1948), 63--66.

\bibitem{hornich1941}
H.~Hornich,
\emph{\"Uber beliebige Teilsummen absolut konvergenter Reihen},
Monatshefte f\"ur Mathematik und Physik \textbf{49} (1941), 316--320.
\href{https://doi.org/10.1007/BF01707309}{doi:10.1007/BF01707309}.

\bibitem{nitecki2013} Z.~Nitecki,
  \href{https://arxiv.org/abs/1106.3779v2}{\emph{Subsum sets: intervals, Cantor
  sets, and Cantorvals}}, arXiv:1106.3779v2 (2013).

\bibitem{fridy1966} John A. Fridy, \href{https://www.fq.math.ca/Scanned/4-3/fridy.pdf}{\emph{Generalized bases for the real numbers}}.
  The Fibonacci Quarterly \textbf{4} (1966), no.~3, 193--201.

\bibitem{crmarickovac2025} T.~Crmari\'{c} and V.~Kova\v{c},
  \href{https://arxiv.org/abs/2504.18712v1}
  {\emph{On the irrationality of certain super-polynomially decaying series}},
  Colloq. Math. \textbf{179} (2025), 55--68,
  doi:\href{https://doi.org/10.4064/cm9628-5-2025}{10.4064/cm9628-5-2025}.
  Theorem locators follow arXiv:2504.18712v1.

\bibitem{kovactao2024} Vjekoslav Kova\v{c} and Terence Tao, \href{https://arxiv.org/abs/2406.17593v4}{\emph{On several irrationality problems for Ahmes series}}.
  Acta Mathematica Hungarica \textbf{175} (2025), 572--608, doi:\href{https://doi.org/10.1007/s10474-025-01528-0}{10.1007/s10474-025-01528-0}; arXiv:\href{https://arxiv.org/abs/2406.17593v4}{2406.17593v4}. Statement and section numbers refer to arXiv version 4.

\bibitem{vandoornkovac2025} Wouter van Doorn and Vjekoslav Kova\v{c}, \href{https://arxiv.org/abs/2509.24971v3}{\emph{Lacunary sequences whose reciprocal sums represent all rational numbers in an interval}}.
  Acta Arithmetica \textbf{223} (2026), 275--295, doi:\href{https://doi.org/10.4064/aa251001-13-1}{10.4064/aa251001-13-1}; arXiv:\href{https://arxiv.org/abs/2509.24971v3}{2509.24971v3}. Statement numbers refer to arXiv version 3.

\bibitem{duverneytachiya} D.~Duverney and Y.~Tachiya,
  \href{https://danielduverney.fr/documents/theorie-des-nombres/DuverneyTachiya190522.pdf}
  {\emph{Refinement of the Chowla--Erd\H{o}s method and linear independence
  of certain Lambert series}},
  Forum Math.\ 31 (2019), no.~6, 1557--1566,
  \href{https://doi.org/10.1515/forum-2018-0299}{DOI}.  Page numbers refer to
  the linked author preprint.

\bibitem{taoteravainen2025} T.~Tao and J.~Ter\"av\"ainen,
  \href{https://arxiv.org/abs/2512.01739v2}{\emph{Quantitative correlations and
  some problems on prime factors of consecutive integers}}, arXiv:2512.01739v2
  (submitted December 2025, revised April 2026).

\bibitem{fan2026comment} Steve Fan,
  \emph{Comment on Erd\H{o}s Problem \#269, thread 269, post 7218} (2026),
  \href{https://www.erdosproblems.com/forum/thread/269#post-7218}{source}. 26 June 2026.

\bibitem{mahler1929} K.~Mahler, \emph{Arithmetische Eigenschaften der L\"osungen
  einer Klasse von Funktionalgleichungen}, Math.\ Ann.\ \textbf{101} (1929),
  342--367.

\bibitem{adamczewskifaverjon2026} Boris Adamczewski and Colin Faverjon,
  \emph{Mahler's method in several variables and finite automata},
  Annals of Mathematics \textbf{204}, no.~2 (2026), 455--533,
  \href{https://doi.org/10.4007/annals.2026.204.2.1}{doi:\nolinkurl{10.4007/annals.2026.204.2.1}}.
  Locators refer to the
  \href{https://faverjon.perso.math.cnrs.fr/AdamczewskiFaverjon_MahlerFiniteAutomata.pdf}{68-page author manuscript}.

\bibitem{boesdarsterdos1981} D.~Boes, R.~Darst and P.~Erd\H{o}s,
  \href{https://www.tandfonline.com/doi/abs/10.1080/00029890.1981.11995266}{\emph{Fat,
  symmetric, irrational Cantor sets}}, Amer.\ Math.\ Monthly \textbf{88} (1981),
  340--341.

\bibitem{rstw2019} A.~Rahm, N.~Solomon, T.~Trauthwein and B.~Weiss,
  \href{https://arxiv.org/abs/1909.01198}{\emph{The distribution of rational
  numbers on Cantor's middle thirds set}}, arXiv:1909.01198.

\bibitem{chowvarjuyu2024} S.~Chow, P.~Varj\'u and H.~Yu,
  \href{https://arxiv.org/abs/2402.18395}{\emph{Counting rationals and
  diophantine approximation in missing-digit Cantor sets}}, arXiv:2402.18395.

\bibitem{plectis68} W.~Cook,
  \href{erdos-68-factorial-denominator-irrationality.pdf}{\emph{Two
  incomparable denominator exclusions for $\sum_{n\ge2}(n!-1)^{-1}$}},
  Erd\H{o}s Problem Notes, 2026.

\bibitem{plectis243} W.~Cook,
  \href{erdos-243-reciprocal-tail-rigidity.pdf}{\emph{Bounded
  increments and rational reciprocal sums}}, Erd\H{o}s Problem Notes, 2026.

\bibitem{plectis249} W.~Cook,
  \href{erdos-249-binary-totient-series.pdf}{\emph{Bases and
  integral relations for the $k$-kernel of Euler's totient}}, Erd\H{o}s Problem
  Notes, 2026.

\bibitem{plectis251} W.~Cook,
  \href{erdos-251-prime-gap-dyadic-series.pdf}{\emph{Sparse
  congruence-preserving perturbations of dyadic series}}, Erd\H{o}s Problem
  Notes, 2026.

\bibitem{plectis257} W.~Cook,
  \href{erdos-257-mersenne-support-subseries.pdf}{\emph{Weighted
  support criteria for reciprocal Mersenne subseries}}, Erd\H{o}s Problem Notes,
  2026.

\bibitem{plectis269} W.~Cook,
  \href{erdos-269-three-prime-running-lcm.pdf}{\emph{No
  finite separable representation at three prime generators}}, Erd\H{o}s Problem
  Notes, 2026.

\bibitem{plectis1049} W.~Cook,
  \href{erdos-1049-rational-base-lambert.pdf}{\emph{Zudilin's
  forms at rational bases and the exact normalised Hankel order}}, Erd\H{o}s
  Problem Notes, 2026.
\end{thebibliography}

\section*{Declaration of generative AI use}
The mathematics and text of this record were produced with large language
model agents under the direction of the author, who is responsible for every
claim.  The checks were these: each proof was written by one agent and
refereed by a second with instructions to break it; the computations were
rerun independently; the external citations were opened on 20 September 2026,
and the one that failed is recorded in Section~\ref{sec:eliminated}.  None of
this is review by a mathematician.
\end{document}
